\scanpage{27}
\mspage{13}
L'hyp. que $\pi_1(\unsure{f'|f})\unsure{=|\simeq}\pi_1(g)$ signifie donc qu'il existe $\alpha\in\pi_1(Y)$, tel que
\[
\pi_1(j')(\sigma)=\inserted{\operatorname{int}(\alpha)}\pi_1(j)(\sigma)
\]
pour tt $\sigma\in\operatorname{Im}(\pi_1(X)\xrightarrow{\pi_1(f)}\pi_1(U))$. Or cette image est un sous-groupe \uline{ouvert} de $\pi_1(U)$ (tjrs vrai pour un morphisme dominant!). \deleted{D'ailleurs} Donc on est ramené à ceci: Soit $U$ ouvert \inserted{$\ne\varnothing$} de $Y$, et $a\in G$, tels que $\tau_aU\subset Y$ [et tels que (désignant par $j,j'$ les \deleted{applications} \inserted{morphismes} $y\longrightarrow y$ et $y\longmapsto y+a$ de $U$ dans $Y$) $\pi_1(j)$ et $\pi_1(j')$ coïncident extérieurement sur un sous-groupe ouvert de $\pi_1(U)$] -- \uline{alors} $j=j'$ i.e. $a=0$.

\margin{\uline{NB} ceci n'est pas plus fort que le corollaire, comme on voit en prenant pour $X$ le rev. étale fini de $U$ correspondant à un sous-groupe ouvert de $\pi_1(U)$.}

Finalement, je m'aperçois que je n'arrive pas à le prouver par voie géométrique -- pour ceci je me rabats sur une méthode \uline{arithmétique}.
\margin{Si -- cf. plus bas!}

\deleted{\unsure{Je ferai} d'ailleurs \deleted{\unsure{faible}} \unsure{généraliser} le th. géométrique; on montre que (quitte à remplacer $X$ par un ouvert non vide connexe, ce qui ne détruit pas les hyp. ni ne modifie la conclusion) OPS que $X$ est une variété}
\margin{26}
